\documentclass[11pt]{article}
\usepackage{amsmath,amssymb,amsthm}
\usepackage[margin=1.1in]{geometry}
\usepackage{hyperref}

\newtheorem{lemma}{Lemma}
\newtheorem{proposition}[lemma]{Proposition}
\newtheorem{corollary}[lemma]{Corollary}
\newtheorem{conjecture}[lemma]{Conjecture}
\theoremstyle{remark}
\newtheorem{remark}[lemma]{Remark}
\newtheorem{question}[lemma]{Question}

\newcommand{\F}{\mathbb{F}}
\newcommand{\R}{\mathbb{R}}
\newcommand{\Ze}{Z_{\mathrm{ev}}}
\newcommand{\rhomax}{\rho}
\DeclareMathOperator{\tr}{tr}

\title{Parity structure of the signed spectral radius,\\
with computational evidence around the Bilu--Linial conjecture}
\author{V.\thanks{Draft note, July 2026. Code and data:
\texttt{signed\_spectra.py}, \texttt{atlas.csv}, \texttt{variance\_onset.csv}.}}
\date{July 2026}

\begin{document}
\maketitle

\begin{abstract}
Let $G$ be a connected graph and $\sigma\colon E(G)\to\{\pm1\}$ a signing of
its adjacency matrix. We record two elementary structural lemmas that
explain, and are suggested by, systematic computation on the minimum
spectral radius $\min_\sigma \rhomax(A_\sigma)$: (i) the spectral radius
depends only on the restriction of the sign character to the
\emph{even-weight} part of the cycle space, so odd cycles influence
$\rhomax$ only in pairs, at second order; (ii) any graph containing $K_4$
admits no signing in which every quadrilateral is unbalanced, an instance
of an $\F_2$-dependency obstruction that motivates a graph parameter
$\beta_L(G)$. We further show that on the hypercube $Q_n$ the local
$\F_2$ condition ``every $4$-cycle is unbalanced'' forces $A_\sigma^2=nI$,
so \emph{every} solution attains $\rhomax=\sqrt n$; Huang's sensitivity
signing is one such solution. Finally we record a decisive negative for a
natural two-sided interlacing approach: the expected characteristic
polynomial of $A_\sigma^2$ is not real-rooted, already for the
quadrilateral, where it equals $(x^2-4x+2)^2+4$. Assembling these
ingredients yields an $\varepsilon$-version of the conjecture in a
dilute regime: every $d$-regular graph whose connected subgraphs at
scale $\log n$ have bounded cycle rank and subcritical
non-backtracking growth admits a signing with
$\rho\le2\sqrt{d-1}\,(1+\delta)^{O(r_0)}(1+o(1))$.
\end{abstract}

\section{Setup}

$G=(V,E)$ is a simple connected graph, $|V|=n$, $|E|=m$, adjacency matrix
$A$. A signing is $\sigma\in\{\pm1\}^E$, encoded by the bit vector
$s\in\F_2^E$ with $\sigma_e=(-1)^{s_e}$; the signed adjacency matrix
$A_\sigma$ has entries $\sigma_{uv}A_{uv}$. We write
$\rhomax(A_\sigma)=\max_i|\lambda_i(A_\sigma)|$. The cycle space
$Z(G)\le\F_2^E$ has dimension $m-n+1$; the \emph{sign} of $z\in Z(G)$
under $\sigma$ is $(-1)^{s\cdot z}$, and a cycle $C$ is \emph{unbalanced}
if $s\cdot\chi_C=1$. Switching preserves the spectrum and acts trivially
on cycle signs; conversely the switching class of $\sigma$ is exactly the
linear functional $z\mapsto s\cdot z$ on $Z(G)$, so there are
$2^{m-n+1}$ switching classes.

For $d$-regular $G$ the conjecture of Bilu and Linial
\cite{bilulinial06} asserts $\min_\sigma\rhomax(A_\sigma)\le
2\sqrt{d-1}$, the spectral radius of the infinite $d$-regular tree.
Marcus, Spielman and Srivastava \cite{mss15} proved the one-sided bound
$\lambda_{\max}(A_\sigma)\le2\sqrt{d-1}$ for some $\sigma$, resolving the
bipartite case; Mohanty, O'Donnell and Paredes \cite{mop20} showed a
random signing attains $\rhomax\le2\sqrt{d-1}+\varepsilon$ when every
ball of radius $r\gg(\log\log n)^2$ contains at most one cycle; and Xu
and Zhang \cite{xuzhang26} recently proved
$\min_\sigma\rhomax(A_\sigma)\le2\sqrt{3(d-1)}$ for all graphs of maximum
degree $d$ via interlacing families of mixed characteristic polynomials.

\section{Even-trace parity}

For a closed walk $W$ of length $\ell$ let $m_e(W)$ be the number of
traversals of $e$, and let $z(W)\in\F_2^E$ be the parity vector
$z(W)_e=m_e(W)\bmod 2$. Since $W$ is closed, every vertex has even total
incidence with the edge multiset of $W$, so the support of $z(W)$ has all
even degrees: $z(W)\in Z(G)$. Moreover
$|z(W)|\equiv\sum_e m_e(W)=\ell\pmod 2$, and the sign of $W$ under
$\sigma$ is $\prod_e\sigma_e^{m_e(W)}=(-1)^{s\cdot z(W)}$.

Let $\Ze=\{z\in Z(G):|z|\text{ even}\}$, the kernel of the linear map
$z\mapsto|z|\bmod2$ on $Z(G)$. If $G$ is bipartite then $\Ze=Z(G)$;
otherwise $\Ze$ has index $2$, hence dimension $m-n$.

\begin{lemma}[Even-trace parity]\label{lem:parity}
For every signing $\sigma$ and every even $\ell$,
\[
\tr\!\big(A_\sigma^{\ell}\big)
=\sum_{z\in \Ze} N_\ell(z)\,(-1)^{s\cdot z},
\]
where $N_\ell(z)\ge0$ is the number of closed $\ell$-walks in $G$ with
parity vector $z$, independent of $\sigma$.
\end{lemma}

\begin{proof}
Group the closed $\ell$-walks by $z(W)$ and use
$\mathrm{sign}(W)=(-1)^{s\cdot z(W)}$. For even $\ell$ only classes with
$|z|$ even occur, by the congruence above.
\end{proof}

\begin{corollary}\label{cor:evenchar}
$\rhomax(A_\sigma)$ is a function of the restriction of the character
$z\mapsto(-1)^{s\cdot z}$ to $\Ze$. In particular
$\rhomax(A_{-\sigma})=\rhomax(A_\sigma)$, and for non-bipartite $G$ the
map from switching classes to $\rhomax$-relevant data is $2$-to-$1$: at
most $2^{m-n}$ values of $\rhomax$ occur among the $2^{m-n+1}$ classes.
\end{corollary}

\begin{proof}
$\rhomax(A_\sigma)=\lim_{k\to\infty}\tr(A_\sigma^{2k})^{1/2k}$ since
$\rhomax^{2k}\le\tr(A_\sigma^{2k})\le n\rhomax^{2k}$, and by
Lemma~\ref{lem:parity} every even trace depends on $s$ only through
$s|_{\Ze}$. Negating $\sigma$ adds $|z|\bmod 2$ to $s\cdot z$, which
vanishes on $\Ze$.
\end{proof}

\begin{corollary}[Interference accounting]\label{cor:interference}
Let $C$ be a cycle of $G$.
If $|C|$ is even, the class $z=\chi_C$ contributes to even traces with
the sign $(-1)^{s\cdot\chi_C}$ of $C$ itself (first order). If $|C|$ is
odd, then $\chi_C\notin\Ze$ and $C$ enters even traces only through
combined classes; the lowest-order such classes are
$z=\chi_{C_1}+\chi_{C_2}$ for pairs of odd cycles, contributing with sign
equal to the \emph{product} of the two cycle signs.
\end{corollary}

Corollary~\ref{cor:interference} identifies the sign structure but not
the magnitudes $N_\ell(z)$, so it does not by itself prove monotonicity
statements. It does, however, explain two robust empirical phenomena
(\S\ref{sec:experiments}): making all short \emph{even} cycles unbalanced
places a deterministic $(-1)$ on their first-order terms, while making
all short \emph{odd} cycles unbalanced makes every odd-pair term
$(+1)(+1)$ or $(-1)(-1)$, i.e.\ sign-definite \emph{positive},
eliminating cancellation; and indeed exhaustive optima place roughly half
of all triangles unbalanced, while forcing all triangles unbalanced
raises $\rhomax$ above the random-signing mean.

\section{An $\F_2$ obstruction and the parameter $\beta_L$}

Fix $L$ and consider the linear system over $\F_2$
\begin{equation}\label{eq:sys}
s\cdot\chi_C=1\qquad\text{for every even cycle $C$ of length}\le L .
\end{equation}

\begin{lemma}[$K_4$ obstruction]\label{lem:k4}
If $G$ contains $K_4$ as a subgraph, then the system \eqref{eq:sys} with
$L\ge4$ is inconsistent: no signing of $G$ makes all quadrilaterals
unbalanced.
\end{lemma}

\begin{proof}
Let $\{a,b,c,d\}$ span a $K_4$ and let $C_1,C_2,C_3$ be its three
$4$-cycles. Every edge of the $K_4$ lies in exactly two of them, so
$\chi_{C_1}+\chi_{C_2}+\chi_{C_3}=0$ in $\F_2^E$, whence
$s\cdot\chi_{C_1}+s\cdot\chi_{C_2}+s\cdot\chi_{C_3}=0$ for every $s$,
while \eqref{eq:sys} demands the sum be $1+1+1=1$.
\end{proof}

More generally, any odd-cardinality dependent family of even short
cycles blocks \eqref{eq:sys}. Define
\[
\beta_L(G)\;=\;\frac{1}{\#\{\text{even cycles}\le L\}}\;
\max_{s\in\F_2^E}\#\{\text{even cycles }C,\ |C|\le L:\ s\cdot\chi_C=1\},
\]
the maximum simultaneously-unbalanceable fraction. Row reduction of
\eqref{eq:sys} certifies upper bounds on $\beta_L$ and produces the
affine solution family of any consistent subsystem; computing $\beta_L$
exactly is MAX-XOR-SAT and NP-hard in general \cite{hastad01}. Our data
(\S\ref{sec:experiments}) suggest $\beta_L$ as the parameter governing
the defect $\min_\sigma\rhomax(A_\sigma)-2\sqrt{d-1}$ on cycle-dense
instances.

\section{Local certificates: the hypercube and $C_n(1,2)$}

\begin{proposition}\label{prop:cube}
On $Q_n$ the system \eqref{eq:sys} with $L=4$ is consistent, and
\emph{every} solution $\sigma$ satisfies $A_\sigma^2=nI$; consequently
$\rhomax(A_\sigma)=\sqrt n$ for the entire affine solution family.
\end{proposition}

\begin{proof}
Consistency is witnessed by Huang's signing \cite{huang19}. Let $\sigma$
be any solution. $Q_n$ is triangle-free, so
$(A_\sigma^2)_{uv}=0$ for adjacent $u,v$. If $u\ne v$ are at distance
$2$ they have exactly two common neighbours $w,w'$, and
$u\,w\,v\,w'$ is a $4$-cycle; unbalancedness gives
$\sigma_{uw}\sigma_{wv}=-\sigma_{uw'}\sigma_{w'v}$, so the two terms of
$(A_\sigma^2)_{uv}$ cancel. Pairs at distance $\ge3$ contribute nothing,
and the diagonal equals the degree $n$.
\end{proof}

This proves the zero variance of $\rhomax$ observed across $200$ sampled
solutions on $Q_5,Q_6$ (all at $\sqrt5,\sqrt6$ exactly, against a Kesten
floor of $4$ and $2\sqrt5$): on $Q_n$ the local $\F_2$ condition is an
exact spectral certificate, far below $2\sqrt{d-1}$. It also isolates
what is special: the certificate needs \emph{every} $2$-path to be
completed by exactly one $4$-cycle. Graphs with this property are
covered by the two-eigenvalue framework of signed covers; the general
question is how much of the cancellation survives when only a
$\beta_L$-fraction of quadrilaterals can be unbalanced.

The circulant $C_n(1,2)$ is a second exactly solvable case, discovered
through the variance experiment.

\begin{proposition}\label{prop:circ}
Let $n\ge7$ be even. On $C_n(1,2)$ the system \eqref{eq:sys} with $L=4$
is consistent and its solutions form exactly four switching classes. The
class containing the signing $\sigma\equiv+1$ on step-$1$ edges
$\{i,i{+}1\}$ and $\sigma_{\{i,i+2\}}=(-1)^i$ on step-$2$ edges has
spectrum
\[
\Big\{\pm2\sqrt{\cos^2\theta_k+\cos^2 2\theta_k}\ :\
\theta_k=\tfrac{2\pi k}{n},\ 0\le k<\tfrac n2\Big\},
\]
so its spectral radius equals $2\sqrt2$ exactly, against the Kesten floor
$2\sqrt3$. For odd $n$ the system is inconsistent, and at most $n-1$
quadrilaterals can be simultaneously unbalanced.
\end{proposition}

\begin{proof}
The quadrilaterals of $C_n(1,2)$ are $Q_i=(i,i{+}1,i{+}3,i{+}2)$, one per
vertex, with sign $a_ia_{i+2}\,b_ib_{i+1}$, where $a_i,b_i$ denote the
signs of $\{i,i{+}1\}$ and $\{i,i{+}2\}$. The displayed signing gives
$a_ia_{i+2}b_ib_{i+1}=(-1)^i(-1)^{i+1}=-1$ for all $i$, so the system is
consistent. If $\sum_{i\in S}\chi_{Q_i}=0$, then comparing coefficients
of the step-$2$ edge $\{i,i{+}2\}\in Q_i\cap Q_{i-1}$ forces $S$ to be
shift-invariant, so $S=\emptyset$ or $S=\mathbb{Z}_n$; the full sum is
indeed zero (each edge lies in exactly two quadrilaterals). Hence the
constraint map on the $2^{\,m-n+1}=2^{\,n+1}$ switching classes has rank
$n-1$, and the solution set is a coset of dimension $2$: four classes.

For the displayed class, write $C_1,C_2$ for the unsigned step-$1$ and
step-$2$ circulants and $D=\operatorname{diag}((-1)^j)$; then
$A_\sigma=C_1+DC_2$, which is symmetric since $D$ commutes with the even
shift $C_2$. On the orthonormal Fourier basis
$f_k(j)=n^{-1/2}e^{2\pi ijk/n}$ we have $C_1f_k=2\cos\theta_k\,f_k$,
$C_2f_k=2\cos2\theta_k\,f_k$, and $Df_k=f_{k+n/2}$. Each plane
$\operatorname{span}\{f_k,f_{k+n/2}\}$ is therefore invariant, carrying
the block
\[
\begin{pmatrix} 2\cos\theta_k & 2\cos2\theta_k\\
2\cos2\theta_k & -2\cos\theta_k\end{pmatrix},
\]
where we used $\cos\theta_{k+n/2}=-\cos\theta_k$ and
$\cos2\theta_{k+n/2}=\cos2\theta_k$. A traceless symmetric $2\times2$
block $\left(\begin{smallmatrix}a&b\\ b&-a\end{smallmatrix}\right)$ has
eigenvalues $\pm\sqrt{a^2+b^2}$, giving the stated spectrum. Setting
$u=\cos^2\theta\in[0,1]$,
$\cos^2\theta+\cos^22\theta=4u^2-3u+1\le2$ with equality iff $u=1$,
attained at $k=0$; hence $\rhomax=2\sqrt2$. For odd $n$, summing all $n$
constraints gives $0$ on the left (each edge appears twice) and
$n\equiv1\pmod2$ on the right.
\end{proof}

\begin{proposition}[Twisted classes]\label{prop:twist}
In the setting of Proposition~\ref{prop:circ}, let $\tau_i$ denote the
sign of the triangle $T_i=(i,i{+}1,i{+}2)$ and $\alpha$ the sign product
around the step-$1$ Hamilton cycle $H$. Then
$\chi_{Q_i}=\chi_{T_i}+\chi_{T_{i+1}}$, so \eqref{eq:sys} is equivalent
to the alternation $\tau_{i+1}=-\tau_i$; the triangles together with $H$
form a basis of the cycle space, so the four solution classes are
coordinatized by $(\tau_0,\alpha)\in\{\pm1\}^2$; and rotation by one
vertex is an automorphism carrying $(\tau_0,\alpha)$ to
$(-\tau_0,\alpha)$. Consequently $\rhomax$ depends on $\alpha$ alone: it
equals $2\sqrt2$ for $\alpha=+1$, and for $\alpha=-1$ it equals
\[
\rho_-(n)\;:=\;2\sqrt{\cos^2(\pi/n)+\cos^2(2\pi/n)}\;<\;2\sqrt2 .
\]
In particular $\min_\sigma\rhomax(A_\sigma)\le\rho_-(n)$ for every even
$n$.
\end{proposition}

\begin{proof}
Each step-$2$ edge $b_i$ lies in $T_i$ only, so
$\sum_{i\in S}\chi_{T_i}=0$ forces $S=\emptyset$; together with the fact
that no combination of triangles has empty step-$2$ support except the
empty one, $\{\chi_{T_0},\dots,\chi_{T_{n-1}},\chi_H\}$ is a basis of
the $(n{+}1)$-dimensional cycle space, giving the coordinates; the
rotation statement is immediate since rotation shifts the alternating
pattern by one. For $\alpha=-1$ it therefore suffices to treat the
representative with $\tau_i=(-1)^i$. Let $R$ be the cyclic shift,
$D=\operatorname{diag}((-1)^j)$, $\varphi=\pi/n$, and set
\[
A(\varphi)\;=\;e^{i\varphi}R+e^{-i\varphi}R^{*}
\;+\;e^{2i\varphi}DR^{2}+e^{-2i\varphi}R^{-2}D ,
\]
a Hermitian matrix supported on $C_n(1,2)$. Its holonomy on the directed
triangle $i\to i{+}1\to i{+}2\to i$ is
$e^{i\varphi}\cdot e^{i\varphi}\cdot(-1)^ie^{-2i\varphi}=(-1)^i$, and on
$H$ it is $e^{in\varphi}=-1$; two Hermitian matrices with common
connected support and equal holonomy on every cycle are conjugate by a
diagonal unitary, so $A(\varphi)$ is isospectral with the real signed
adjacency matrix of the class $(+1,-1)$. Note the step-$2$ phase is
forced to be $2\varphi$: this is exactly the co-shift that keeps the
triangle holonomies real, and it is the step omitted by a naive
one-parameter Bloch argument. As in Proposition~\ref{prop:circ}, the
planes $\operatorname{span}\{f_k,f_{k+n/2}\}$ are invariant, with blocks
\[
\begin{pmatrix}
2\cos(\theta_k{+}\varphi)&2\cos(2\theta_k{+}2\varphi)\\
2\cos(2\theta_k{+}2\varphi)&-2\cos(\theta_k{+}\varphi)
\end{pmatrix},
\]
hence spectrum $\pm2\sqrt{g(\theta_k+\pi/n)}$ with
$g(\theta)=\cos^2\theta+\cos^22\theta$ and shifted momenta
$\theta_k+\pi/n\in\{(2k{+}1)\pi/n\}$, which avoid $0$ and $\pi$. Writing
$g=4u^2-3u+1$ with $u=\cos^2\theta$, the maximum over the shifted
lattice is attained at the momentum nearest $0$, namely $\pi/n$, giving
$\rho_-(n)$.
\end{proof}

\begin{conjecture}\label{conj:circ}
For every even $n\ge8$,
$\min_\sigma\rhomax(A_\sigma)=\rho_-(n)$ on $C_n(1,2)$: the $\alpha=-1$
twisted class is the global optimizer. This is verified by exhaustive
enumeration of all $2^{n+1}$ switching classes for
$n\in\{8,10,12,14,16,18\}$, with agreement to $10^{-9}$. The lower
bound over all classes is a flux-minimization statement in the spirit of
Lieb's flux-phase theorem \cite{lieb94}: among all $\F_2$ flux
assignments to the triangle--Hamilton cycle basis, alternating triangle
flux with anti-periodic step-$1$ holonomy minimizes the spectral radius.
On this family the quadrilateral rule thus contains the (conjecturally
exact) global optimizer, while random signings concentrate \emph{above}
the Kesten floor $2\sqrt3$ from $n\approx480$.
\end{conjecture}

\section{Failure of two-sided interlacing for $A_\sigma^2$}

By Godsil--Gutman \cite{godsilgutman81},
$\mathbb{E}_\sigma\det(xI-A_\sigma)=\mu_G(x)$, the matching polynomial,
which is real-rooted with roots in $[-2\sqrt{d-1},2\sqrt{d-1}]$
\cite{heilmannlieb72}; this underlies \cite{mss15}. A natural two-sided
analogue would take
$q_G(x):=\mathbb{E}_\sigma\det\!\big(xI-A_\sigma^2\big)$ and hope for
real-rootedness plus an interlacing family, which would give some
$\sigma$ with $\rhomax(A_\sigma)^2\le\operatorname{maxroot}(q_G)$.

\begin{remark}\label{rem:qG}
$q_G$ is not real-rooted in general. For $G=C_4$ the two switching
classes have $A_\sigma^2$-spectra $\{4,4,0,0\}$ and $\{2,2,2,2\}$, so
\[
q_{C_4}(x)=\tfrac12\big[x^2(x-4)^2+(x-2)^4\big]
=x^4-8x^3+20x^2-16x+8=(x^2-4x+2)^2+4,
\]
which has no real roots at all. Exhaustive computation over all
connected regular graphs on at most $7$ vertices, plus Petersen, $Q_3$
and $K_{3,3}$, shows real-rootedness fails for every instance containing
an even cycle; it holds trivially for odd cycles (where
$A_\sigma^2$ is signing-independent) and, curiously, for $K_4$. In all
computed cases every complex root satisfied
$\operatorname{Re}\le4(d-1)$.
\end{remark}

The even-subgraph corrections that cancel in
$\mathbb{E}\det(xI-A_\sigma)$ survive squaring; this is the concrete
obstruction to extending the interlacing method two-sidedly over the
plain signing family, and delimits where any improvement of
\cite{xuzhang26} must diverge from it.

\section{Computational evidence}\label{sec:experiments}

We summarize the computations (code in \texttt{signed\_spectra.py};
switching classes enumerated exactly as co-tree sign patterns).

\emph{Defect atlas.} For every connected regular graph on at most $7$
vertices and for Petersen, $Q_3$, $Q_4$, Heawood, M\"obius--Kantor and
circulants $C_n(1,2)$, the exact minimum of $\rhomax$ over all switching
classes is at most $2\sqrt{d-1}$, usually strictly below
($K_6$: $\sqrt5\approx2.236$ against $3.464$, attained by a conference
signing; octahedron: $2.0$). At these sizes $75$--$97\%$ of switching
classes lie strictly below the floor. The exhaustive optimum on $Q_4$ is
Huang's signing. Across all optima, quadrilaterals are unbalanced at
rate $\approx1$, triangles at rate $\approx\tfrac12$, consistent with
Corollary~\ref{cor:interference}.

\emph{The even-cycle rule at scale.} Solving \eqref{eq:sys} and sampling
the affine family: on $Q_5,Q_6$ all constraints are satisfiable and
$\rhomax=\sqrt n$ exactly (Proposition~\ref{prop:cube}); on $C_n(1,2)$ with $n$ even the
quadrilateral system is fully satisfiable and the family sits at
$\rhomax\le2\sqrt2\approx2.83$ against a floor of $3.46$
(Proposition~\ref{prop:circ}), matching local search at negligible cost. Forcing triangles
unbalanced instead \emph{raises} $\rhomax$ to $3.95$, above the
random-signing mean $3.44$. On a locally tree-like control (random
cubic, $n=60$) all methods sit at the floor, as theory requires. On
$K_{12}$ the system is $94\%$ inconsistent (Lemma~\ref{lem:k4}) and the
rule collapses; dense cliques belong to Seidel-matrix territory.

\emph{Variance onset.} On symmetry-free instances (random $2$-lift
towers of $Q_3$ and $K_4$, planted-quadrilateral random $4$-regular
graphs) the spread of $\rhomax$ across the solution family of
\eqref{eq:sys} is nonzero from $n\approx32$ on and its standard
deviation stays roughly constant ($\approx1.5\times10^{-2}$ on towers)
while the random-signing standard deviation decays; the minimum over
$200$ sampled solutions remains below the Kesten floor at every size
tested (at $n=1024$: family standard deviation
$\approx1.3\times10^{-2}$, now $2.4$--$2.5$ times the random-signing
deviation; minimum over $200$ tails at floor $-0.02$, tracking the
Alon--Boppana-forced profile), while the zero-tail particular solution
drifts slowly above it (floor $+0.15$ at $n=1024$). Long-cycle terms
are therefore not negligible on generic graphs, and any trace-method
error analysis must be quantitative in them; the transitive flatness of
$Q_n$ and circulants is explained by Proposition~\ref{prop:cube} and by
symmetry respectively.

\section{The kernel-averaged trace}\label{sec:kernel}

Fix a consistent system \eqref{eq:sys}, let
$W=\operatorname{span}\{\chi_C\}\subseteq\Ze$ be the span of its
constraint cycles, and let $\mathcal F$ be its family of switching
classes.

\begin{proposition}[Master identity]\label{prop:master}
The map $\pi\colon W\to\F_2$,
$\pi\big(\sum_{C\in S}\chi_C\big)=|S|\bmod2$, is a well-defined linear
form, and for every even $\ell$,
\[
\mathbb E_{\sigma\in\mathcal F}\,\tr\!\big(A_\sigma^{\ell}\big)
\;=\;\sum_{z\in W}(-1)^{\pi(z)}\,N_\ell(z).
\]
\end{proposition}

\begin{proof}
Consistency means every dependency $\sum_{C\in S}\chi_C=0$ has $|S|$
even, so $\pi$ is well defined; linearity follows from
$|S|+|S'|\equiv|S\triangle S'|\pmod2$. Identify switching classes with
$\operatorname{Hom}(Z(G),\F_2)$; then
$\mathcal F=\varphi_0+\operatorname{Ann}(W)$ and $\varphi|_W=\pi$ for
every $\varphi\in\mathcal F$. For $z\in Z(G)$ the average
$\mathbb E_{\psi\in\operatorname{Ann}(W)}(-1)^{\psi(z)}$ equals $1$ if
$z\in(\operatorname{Ann}W)^{\perp}=W$ and $0$ otherwise. Averaging the
expansion of Lemma~\ref{lem:parity} term by term gives the claim; note
$W\subseteq\Ze$ since the constraint cycles are even.
\end{proof}

Walks whose wrap parity touches any cycle outside $W$ are annihilated;
a single constraint cycle enters with $-1$, a pair with $+1$, and so on.

\begin{corollary}\label{cor:certificate}
For every $k$,
\[
\min_{\sigma\in\mathcal F}\rhomax(A_\sigma)\;\le\;
R_{\mathcal F}(k):=\Big(\sum_{z\in W}(-1)^{\pi(z)}N_{2k}(z)\Big)^{1/2k}.
\]
By contrast the average over all signings retains only $z=0$:
$\mathbb E_\sigma\tr(A_\sigma^{2k})=N_{2k}(0)\ge n\,t_{2k}(d)$, where
$t_{2k}(d)$ counts closed walks at the root of the infinite $d$-regular
tree; indeed every such tree walk crosses each tree edge an even number
of times and projects, injectively by unique lifting, to a closed walk
in $G$ with all edge multiplicities even. Hence the uniform average can
never certify a spectral radius below the Kesten profile, while
$R_{\mathcal F}(k)$ can, if and only if the signed sum over
$W\setminus\{0\}$ is negative enough to consume both the even-wrap
excess $N_{2k}(0)-n\,t_{2k}(d)$ and part of the tree term.
\end{corollary}

\begin{corollary}[Confined mass]\label{cor:balanced}
The homogeneous family $\operatorname{Ann}(W)$ --- all constraint cycles
\emph{balanced} --- measures the total parity-confined walk mass:
\[
\mathbb E_{\psi\in\operatorname{Ann}(W)}\,\tr\!\big(A_\psi^{2k}\big)
=\sum_{z\in W}N_{2k}(z).
\]
This is the same annihilator computation with $\varphi_0=0$, so every
$z\in W$ enters with sign $+1$; it makes the confined mass, and hence
Question~\ref{q:confined} below, directly measurable by sampling
eigenvalues.
\end{corollary}

Numerically the mechanism is already decisive at $k=40$ with $400$
sampled family members ($n=128$). On a random $2$-lift tower of $K_4$
($d=3$, Kesten floor $2\sqrt2\approx2.828$):
$R_{\mathcal F}(40)\approx2.767$ against
$R_{\mathrm{all}}(40)\approx2.846$ and tree rate $2.839$ --- the
kernel-averaged trace alone certifies a sub-Kesten signing on a
symmetry-free instance, which Corollary~\ref{cor:certificate} shows the
uniform average can never do. On a planted-quadrilateral $4$-regular
graph ($\beta_4=37/41$, floor $3.464$):
$R_{\mathcal F}(40)\approx3.327$. On $C_{60}(1,2)$ the family average is
dominated by its $2\sqrt2$-class, $R_{\mathcal F}(40)\approx2.872$,
while $R_{\mathrm{all}}$ stalls near $3.68$, above the tree rate, as the
theory demands.

Measuring all three averages against the tree profile $n\,t_{2k}(d)$
separates the confined masses. On random $2$-lift towers the number of
short cycles stays bounded as $n$ grows, so the confined mass dilutes
and the instances drift toward the bicycle-free regime; the demanding
test is constant local density. On planted-quadrilateral $4$-regular
graphs at fixed density $\tfrac12$ quadrilateral per vertex,
$n\in\{128,256,512\}$, $k=30$: the ordering
$R_{\mathcal F}<R_{\mathrm{all}}\approx R_{\mathrm{tree}}<
R_{\mathrm{bal}}$ holds throughout; the confined-rate excess
$R_{\mathrm{bal}}-R_{\mathrm{all}}\approx0.12$--$0.15$ is
$n$-independent; the confined-mass ratio
$\mathbb E_{\mathrm{bal}}/\mathbb E_{\mathrm{all}}$ remains $O(1)$ in
$n$ at fixed $k$; and
$\mathbb E_{\mathcal F}\tr(A^{2k})/(n\,t_{2k})\approx0.08$--$0.11$,
below one with an $n$-stable margin --- the parity-corrected average
beats the tree profile itself. These are precisely the behaviours the
counting question below predicts, with rate excess $\delta\approx0.04$
at this density.

\section{A kernel-averaged $L$-function}\label{sec:zeta}

By Bass's theorem \cite{bass92}, for any signing
$\det(I-uB_\sigma)=\prod_{p}\big(1-\sigma(p)u^{|p|}\big)$, the product
running over primes (primitive cyclic classes of closed non-backtracking
walks), where $\sigma(p)=(-1)^{\varphi(z(p))}$ and $z(p)\in Z(G)$ is the
edge-multiplicity parity vector of $p$. Averaging over the family
diagonalizes this product.

\begin{proposition}[Kernel-averaged $L$-identity]\label{prop:Lident}
With $W$, $\pi$, $\mathcal F$ as above, $\varepsilon_p=(-1)^{\pi(z(p))}$,
and $B$ the \emph{unsigned} non-backtracking operator, for
$|u|<\tfrac1{d-1}$,
\[
\mathbb E_{\sigma\in\mathcal F}\,\log\det(I-uB_\sigma)
=\tfrac12\log\det(I-u^2B)
-\tfrac12\!\!\sum_{p\,:\,z(p)\in W}\!\!
\log\frac{1+\varepsilon_pu^{|p|}}{1-\varepsilon_pu^{|p|}}\,.
\]
\end{proposition}

\begin{proof}
As in Proposition~\ref{prop:master}, $\varphi(z)$ is frozen at $\pi(z)$
over $\mathcal F$ when $z\in W$ and is a fair coin otherwise. Hence
$\mathbb E\log(1-\sigma(p)u^{|p|})$ equals
$\log(1-\varepsilon_pu^{|p|})$ for parity-confined primes and
$\tfrac12\log(1-u^{2|p|})$ for the rest. Collecting the coin factor over
\emph{all} primes gives
$\tfrac12\log\prod_p(1-(u^2)^{|p|})=\tfrac12\log\det(I-u^2B)$, and the
per-prime correction for confined primes is
$\log(1-\varepsilon u^{\ell})-\tfrac12\log(1-u^{2\ell})
=\tfrac12\log\frac{1-\varepsilon u^{\ell}}{1+\varepsilon u^{\ell}}$.
Absolute convergence for $|u|(d-1)<1$ follows from
$\#\{p:|p|=\ell\}\le\tfrac{2m}{\ell}(d-1)^{\ell}$.
\end{proof}

Both sides transfer to the adjacency side member-wise via
$\det(I-uB_\sigma)=(1-u^2)^{m-n}\det(I-uA_\sigma+u^2(d-1)I)$. The
identity was verified exactly on the prism ($\mathcal F$ of size $2$,
$W$ spanned by the three quadrilaterals): the coefficient form
$\mathbb E_{\mathcal F}\tr B_\sigma^k$ matches the prime decomposition
for all $k\le14$ to machine precision, and the displayed identity holds
to $8$ digits at $u=0.2$ under truncation at $|p|\le14$.

The structural content: $\tfrac12\log\det(I-u^2B)$ is analytic exactly
in $|u|<(d-1)^{-1/2}$, the Ramanujan radius --- \emph{primes whose
parity escapes $W$ contribute the Ramanujan rate automatically, and the
entire deviation of the averaged $L$-function from that rate is carried
by parity-confined primes.} The sign problem has become a counting
problem.

\begin{question}[Confined-prime counting]\label{q:confined}
For which graph sequences does
$\#\{p:|p|=\ell,\ z(p)\in W\}\le C\,(\sqrt{d-1}+\delta)^{\ell}$ hold?
Parity confinement means the odd-multiplicity edge set of $p$ is a sum
of short even cycles; the extreme case $z(p)=0$ (totally even primes)
counts non-backtracking closed walks retracing every edge evenly, which
grow at rate $\sqrt{d-1}$, and each nonzero $z\in W$ should contribute a
comparable rate with an exponential penalty in $|z|$. Caution from the
prism: on graphs so small that $W$ nearly fills the cycle space,
confinement is no restriction (measured rate $\approx1.84$ against
$\sqrt2\approx1.41$ at $d=3$); the question is asymptotic, with
$\dim W$ per ball bounded while $\dim Z(G)\to\infty$.
\end{question}

\section{The $z=0$ slice: totally even walks}\label{sec:even}

Sub-question of Question~\ref{q:confined} at $z=0$: count closed
non-backtracking walks all of whose edge multiplicities are even
(``totally even'' walks). The lower bound holds at the conjectured rate
unconditionally; the upper bound fails without a hypothesis on dense
cores, and the failure mechanism is instructive.

\begin{proposition}[Slice identity and doubling]\label{prop:slice}
Let $\mathrm{Ev}_\ell$ denote the number of totally even closed
non-backtracking walks of length $\ell$ in $G$. Then
$\mathrm{Ev}_\ell=\mathbb E_{\sigma}\,\tr(B_\sigma^{\ell})$, the average
over all signings; and the doubling map $w\mapsto w^2$ injects closed
non-backtracking walks of length $\ell/2$ into totally even walks of
length $\ell$, so
$\mathrm{Ev}_\ell\ge\tr(B^{\ell/2})$ and
$\liminf_\ell \mathrm{Ev}_\ell^{1/\ell}\ge\rho(B)^{1/2}=\sqrt{d-1}$.
\end{proposition}

\begin{proof}
$\tr(B_\sigma^\ell)$ is the signed count of closed non-backtracking
walks, the sign being $\prod_e\sigma_e^{m_e(w)}$; averaging kills every
walk with an odd multiplicity. If $w$ is cyclically non-backtracking
then so is $w^2$ (the seam repeats a legal transition of $w$), its
multiplicities double, and $w$ is recovered as the first half.
\end{proof}

\begin{proposition}[Dense trapping]\label{prop:trap}
Let $d\ge4$ and suppose $G$ contains $K_d$ as a subgraph (realizable in
connected $d$-regular graphs: two copies of $K_d$ joined by a perfect
matching). Then there are $c_d>0$ and $\ell_0(d)$, independent of $G$
and $n$, with $\mathrm{Ev}_\ell(G)\ge c_d\,(d-2)^{\ell}$ for all even
$\ell\ge\ell_0$. Since $d-2>\sqrt{d-1}$ for every $d\ge4$, the confined
mass exceeds the profile $n(\sqrt{d-1}\,)^{\ell}$ throughout
$\ell\ge C_d\log n$.
\end{proposition}

\begin{proof}
Totally even walks of the $K_d$ copy are totally even walks of $G$, and
by Proposition~\ref{prop:slice} their number is
$2^{-|E(K_d)|}\sum_{\sigma\in\{\pm1\}^{E(K_d)}}\tr(B_{K_d,\sigma}^\ell)$.
The unsigned $B_{K_d}$ is irreducible and aperiodic with Perron value
$d-2$; its remaining eigenvalues have modulus at most
$\max(1,\sqrt{d-2})<d-2$ (Ihara--Bass from
$\operatorname{spec}A(K_d)=\{d-1,(-1)^{d-1}\}$). By the equality case of
Wielandt's theorem, a signed version satisfies
$\rho(B_{K_d,\sigma})=\rho(B_{K_d})$ if and only if
$B_{K_d,\sigma}=\pm D B_{K_d}D^{-1}$ for a diagonal $D$ with
$|D|=I$, i.e.\ iff $\sigma$ is switching-equivalent to the all-plus or
all-minus signing. Hence $\rho_2:=\max$ over the remaining classes is
strictly below $d-2$, the two trivial classes contribute
$\tr B^\ell+(-1)^\ell\tr B^\ell=2(d-2)^\ell(1+o(1))$ at even $\ell$, and
the class average is at least
$2^{-r}\big[2(d-2)^\ell(1-o(1))-2^{r}\,d(d-1)\rho_2^{\ell}\big]$ with
$r=\binom d2-d+1$.
\end{proof}

\begin{remark}[Necessity of a core hypothesis]\label{rem:core}
Since $0\in W$ always, Proposition~\ref{prop:trap} shows the bound of
Question~\ref{q:confined} cannot hold unconditionally once
$\ell\ge C_d\log n$: it forces the exclusion of $K_d$ --- for $d=4$ the
same $K_4$ as Lemma~\ref{lem:k4}, which thus obstructs the program
twice --- and, more generally, of any subgraph $H$ whose non-trivial
switching classes carry non-backtracking spectral radius above
$(1+\delta)\sqrt{d-1}$. Question~\ref{q:confined} should accordingly be
read under a \emph{subcritical-cores} hypothesis: every connected
subgraph on at most $\ell_0$ edges has all non-$\pm$-trivial classes
with $\rho(B_{H,\sigma})\le(1+\delta)\sqrt{d-1}$. The data suggest this
is close to sufficient: on $K_4$ ($d=3$) and on the Petersen graph,
\emph{every} non-$\pm$-trivial switching class has
$\rho(B_\sigma)=\sqrt2=\sqrt{d-1}$ exactly, i.e.\ every nontrivial
signing of these graphs is exactly Ramanujan; while on two joined
$K_4$'s ($d=4$) the largest non-trivial class reaches
$\rho\approx2.445$ and the trivial-class-subtracted even-walk rate runs
at $\approx2.1$--$2.2$ against $\sqrt3\approx1.73$, confirming the
trapping. The slice identity itself was verified against brute-force
walk enumeration on the prism through $\ell=10$.
\end{remark}

The upper bound matching Proposition~\ref{prop:slice} holds in the
dilute regime. Say $G$ is \emph{$(r_0,\delta)$-subcritical at scale
$\ell$} if every connected subgraph $H$ of $G$ with minimum degree
$\ge2$ and at most $\ell$ edges satisfies: (i) the cycle rank of $H$
is at most $r_0$; and (ii) the number $c_H(t)$ of closed
non-backtracking walks of length $t$ in $H$ obeys
$c_H(t)\le 2|E(H)|\,(t+1)^{3r_0}(1+\delta)^{t}$ for all $t\le8\ell$.
The notion is monotone --- subcriticality at scale $\ell$ implies it
at every smaller scale --- and all subsequent hypotheses are stated
through this single definition.
Condition (ii) holds, for instance, whenever every thread of $H$ (a
maximal path through degree-$2$ vertices) has length at least
$\lambda$, with $1+\delta=(2r_0-1)^{1/\lambda}$: closed walks in $H$
project to reduced cyclic words in a free group of rank $\le r_0$ whose
generators cost at least $\lambda$ steps each.

\begin{lemma}[Non-backtracking reachability]\label{lem:reach}
Let $H$ be connected with minimum degree $\ge2$, $s$ edges, and not a
cycle. Between any two directed edges of $H$ there is a
non-backtracking path of length at most $4s$.
\end{lemma}

\begin{proof}
Since $H$ has minimum degree $2$ and is not a cycle, some vertex has
degree $\ge3$, so the cycle rank is at least $2$ and $H$ contains two
independent cycles; connectivity yields either a theta subgraph or a
dumbbell $C$---$P$---$C'$ (two cycles joined by a path, possibly of
length zero). Either configuration reverses direction of travel:
arriving at the attachment vertex $x$ of $C$ along the last edge of
$P$, the walk may enter $C$ --- the first $C$-edge at $x$ is not the
reverse of $P$'s last edge --- traverse $C$ once, and exit into the
reverse of $P$'s last edge, which is legal since that edge differs
from the $C$-edge just used; the walk now travels $P$ in the opposite
orientation. Routing from $e$ to $f$ proceeds in at most four phases:
reach the reversing configuration, reverse if the current orientation
disagrees with a fixed non-backtracking route toward $f$, follow that
route, and align with $f$, using the configuration once more if
needed. A non-backtracking continuation always exists at minimum
degree $2$, and each phase can be chosen to traverse each thread of
$H$ at most once, giving total length at most $4s$.
\end{proof}

\begin{proposition}[Even-walk upper bound, dilute regime]\label{prop:upper}
If $G$ is $d$-regular and $(r_0,\delta)$-subcritical at scale $\ell$,
then
\[
\mathrm{Ev}_\ell(G)\;\le\;
nd\,(C\ell)^{\,3(r_0+1)^2}\,
\Big((1+\delta)^{2(r_0+2)}\sqrt{d-1}\Big)^{\ell}
\]
for an absolute constant $C$. Combined with
Proposition~\ref{prop:slice}, in this regime
$\mathrm{Ev}_\ell=\big(\sqrt{d-1}\big)^{\ell(1+o(1))}$: the $z=0$ case
of Question~\ref{q:confined} is settled.
\end{proposition}

\begin{proof}
Fix the starting directed edge ($\le nd$ choices) and expose the walk
step by step; let $S_i$ be the subgraph of edges used through step $i$
and $S=S_{\ell}$ the support. Since multiplicities are even,
$s:=|E(S)|\le\ell/2$; since a degree-$1$ vertex of $S$ would force a
backtrack, $S$ has minimum degree $2$; by (i) its cycle rank $r$ is at
most $r_0$.

Call a step \emph{fresh} if it traverses an edge not in $S_{i-1}$,
\emph{stale} otherwise. The key structural fact: after a fresh step
into a previously unvisited vertex $v$, the only used edge at $v$ is
the arrival edge, which is excluded by non-backtracking --- so every
continuation is again fresh. Hence maximal fresh runs continue until
they land on an already-visited vertex, i.e.\ each maximal fresh run is
an open ear attached to the current support, and the number of fresh
runs equals the number of ears, which is exactly the cycle rank:
\emph{the walk performs precisely $r\le r_0$ fresh runs} (the first
builds the initial cycle), interleaved with at most $r+1$ maximal stale
segments.

Encode the walk by: the start edge; the set of run-start times, at most
$\binom{\ell}{r}\le\ell^{\,r_0}$ choices (run \emph{ends} are visible
to a decoder who knows $G$ and the visited set: a run ends exactly when
its landing vertex is old); the fresh-edge choices, at most $(d-1)^{s}$
in total; and the stale segments. A stale segment of length $t$ is a
non-backtracking path inside $S$, and by Lemma~\ref{lem:reach} any such
path prolongs to a closed walk of length $\le t+4s$, so their number is
at most $2s\cdot c_S(t+4s)\le 4s^2(t+4s+1)^{3r_0}(1+\delta)^{t+4s}$ by
(ii) (if $S$ is a cycle the count is $1$). Multiplying over the
$\le r_0+1$ segments with total stale length $\ell-s$ and collecting,
\[
\mathrm{Ev}_\ell\le nd\,\ell^{\,r_0}\!\!\sum_{s\le\ell/2}\!
\big(4s^2(4\ell)^{3r_0}\big)^{r_0+1}
(d-1)^{s}(1+\delta)^{(\ell-s)+4s(r_0+1)} .
\]
Every factor of the summand is nondecreasing in $s$, so it is maximized
at $s=\ell/2$, where
$(d-1)^{\ell/2}(1+\delta)^{\ell/2+2\ell(r_0+1)}\le
\big((1+\delta)^{2(r_0+2)}\sqrt{d-1}\big)^{\ell}$; the polynomial
prefactors collect into $(C\ell)^{3(r_0+1)^2}$.
\end{proof}

\begin{remark}[Scope]\label{rem:scope}
Random $2$-lift towers are $(O(1),o(1))$-subcritical at scale
$c\log n$ with high probability --- balls of that radius have bounded
cycle rank and long threads --- and the measured
$\mathbb E_\sigma[\tr B_\sigma^{\ell}]^{1/\ell}\approx1.53$ at
$\ell=30$, $n=128$, $d=3$ against $\sqrt2$ is consistent with the
proposition's $(1+\delta)$-corrected rate. Constant cycle density
(quadrilaterals proportional to $n$) violates hypothesis (i): supports
with $s$ edges can have rank $\Theta(\eta s)$. The run-schedule cost
survives this --- $\binom{\ell}{r}\le(e\ell/r)^{r}=(1+o_\eta(1))^{\ell}$
for $r\le\eta\ell$ --- so the sole obstruction to the constant-density
extension is the stale-path entropy bound for unbounded-rank supports,
which is precisely the quantitative content of the subcritical-cores
hypothesis of Remark~\ref{rem:core}. That extension, together with the
$z\ne0$ insertion accounting, is what remains of
Question~\ref{q:confined}.
\end{remark}

\begin{remark}[Toward constant density: the critical gadget]\label{rem:gadget}
At $d=3$ the candidate obstruction to the unbounded-rank entropy bound
is a quadrilateral with exits at two opposite vertices: oscillation
between its branch points along the two length-$2$ arcs would give
path growth $2^{2/4}=\sqrt2=\sqrt{d-1}$, exactly critical.
Non-backtracking forbids it: arriving along one short arc, a walk may
not reverse into it, so bouncing between the two short arcs is
deterministic, and branching entropy is carried only by the long arm.
Exactly: on the theta graph $\Theta(2,2,\lambda)$ \emph{every}
switching class satisfies $\rho(B_\sigma)\le\rho(B)=1+O(1/\lambda)$
(measured $1.348, 1.268, 1.222, 1.190, 1.150$ for
$\lambda=3,5,7,9,13$, all nontrivial classes strictly below), and a
ring of four such quadrilaterals separated by $\lambda$-threads
behaves identically ($\rho\le1.25$ at $\lambda=6$, rank $5$). The
general principle: if every ball of radius $R$ contains at most one
cycle, then at any branch vertex of any support at most one
continuation returns within $2R$ steps --- two short returns would
exhibit two cycles in one ball. The remaining entropy bound therefore
takes the quantified form: paths in any support grow at rate
$1+O(\log d/R)$; ranks obey $r\le C\ell/R$, keeping the run-schedule
cost $\binom{\ell}{r}$ at $(1+o(1))^{\ell}$; and the per-segment
polynomial factors force $R\ge C\log\log n/\delta$ --- the scale at
which the bicycle-free hypothesis of \cite{mop20} reappears, now in
service of the parity family rather than the uniform measure.
\end{remark}

The entropy half of that program is now a lemma.

\begin{lemma}[Window path bound]\label{lem:window}
Suppose every ball of radius $R$ in $G$ contains at most one cycle.
Then for every subgraph $S\subseteq G$ with $s$ edges, every directed
edge of $S$, and every $t\ge1$, the number of non-backtracking paths
of length $t$ in $S$ from that edge is at most
$\big(2s(R+3)\big)^{\lceil t/R\rceil}$; in particular the growth rate
is $\exp\!\big(O(\log(sR)/R)\big)$.
\end{lemma}

\begin{proof}
Cut the path into $\lceil t/R\rceil$ windows of at most $R$
consecutive steps. A window starting at the vertex $w$ stays inside
$B_R(w)$, so its segment is a non-backtracking path in
$H_w:=S\cap B_R(w)$, a graph of cycle rank at most $1$. In a tree a
non-backtracking walk never revisits a vertex --- a revisit would have
to undo an edge crossing, and the first undoing step is a backtrack
--- so it is the unique reduced path between its endpoints. In a
unicyclic graph with cycle $C$, each component of the complement of
$E(C)$ is a tree attached to $C$ at a single vertex; once a
non-backtracking walk enters such a tree it can never return to the
attachment vertex, and while on $C$ its rotational direction cannot
reverse. A non-backtracking path therefore consists of a unique
reduced tree path to an attachment vertex, unidirectional motion
around $C$ with some winding number $w\le t/3$, and a unique reduced
tree path to its final edge: it is determined by its endpoint directed
edges, the direction, and $w$, giving at most $2(1+t/3)$ paths per
endpoint pair (numerically the sharper entrywise bound
$(B_H^{\,t})_{ef}\le1+t/\mathrm{girth}$ holds with equality). Each
window therefore admits at most $4s(R/3+2)\le2s(R+3)$
continuations given its starting edge, using $R\ge3$.
\end{proof}

The rate is sharp in the exponent. The \emph{tree-burst} gadget ---
two depth-$D$ binary trees with roots and matched leaves joined by
$\lambda$-threads, every cycle of length $\ge2\lambda+4$, hence
bicycle-free at radius $\lambda$ --- has measured $\rho(B)$ tracking
the burst rate $2^{D/(D+\lambda)}$: $(1.29,1.19,1.34,1.22,1.25)$
against $(1.35,1.21,1.41,1.26,1.31)$ for
$(D,\lambda)=(3,4),(3,8),(4,4),(4,8),(5,8)$. Since $D$ can reach
$\log_2(s/\lambda)$, path growth $\exp(c\log s/R)$ is realized, and
$R\ge C\log\ell/\delta$ is \emph{necessary} for a $(1+\delta)^t$
bound, matching Lemma~\ref{lem:window}.

\begin{lemma}[Rank of bicycle-free graphs]\label{lem:rank}
There is an absolute constant $C$ such that every graph $H$ with
$s\ge R\ge8$ edges and minimum degree $\ge2$ that is bicycle-free at
radius $R$ satisfies
$\operatorname{rank}(H)\le1+Cs\log s/R$.
\end{lemma}

\begin{proof}
Call a cycle \emph{short} if its length is at most $2R$. Two short
cycles cannot share a vertex, since both would lie in that vertex's
$R$-ball; hence short cycles are pairwise vertex- and edge-disjoint,
and there are at most $s/R$ of them with length in $(R,2R]$. Two
disjoint short cycles of length $\le R$ at distance $\rho$ have
$\rho>R/2$: otherwise, taking $v$ on one nearest the other, every
vertex of both lies within $\rho+R/2\le R$ of $v$, giving a rank-$2$
ball. Their $R/4$-balls are then disjoint, and each contains at least
$R/4$ edges by minimum degree $2$, so there are at most $4s/R$ of
them. Delete one edge from each short cycle: $D\le5s/R$ deletions
destroy every short cycle and create none, so the result $H'$ has
girth exceeding $2R$. Pass to the $2$-core of each component of $H'$;
this changes neither rank nor girth. By the Moore bound for irregular
graphs \cite{ahl02}, a graph with minimum degree $\ge2$, girth
$\ge2R+1$, and average degree $\bar d_i$ on $n_i$ vertices satisfies
$n_i\ge(\bar d_i-1)^{R}$ for either girth parity, so
$\bar d_i-1\le n_i^{1/R}$, whence
$\operatorname{rank}(H')\le\sum_i\big[\tfrac{n_i}2\big(n_i^{1/R}-1\big)+1\big]$.
If $\log s\le R$ this is at most $Cs\log s/R+D+1$, using
$e^{x}-1\le2x$ on $[0,1]$, $\sum_in_i\le s$, and
$\#\mathrm{components}\le D+1$; if $\log s>R$ the claimed bound
exceeds $s$ and is trivial. Finally
$\operatorname{rank}(H)\le\operatorname{rank}(H')+D$.
\end{proof}

The $\log s$ factor is necessary, and the earlier guess
$\operatorname{rank}\le1+Cs/R$ is false for every constant $C$: a
graph of girth $\ge2R+2$ has every $R$-ball a tree, hence is
bicycle-free at radius $R$, while its rank density is $1-2/\bar d$.
The Tutte--Coxeter $(3,8)$-cage is bicycle-free at radius $3$ with
$(\operatorname{rank}-1)/s=\tfrac13$, the girth-$12$ Tutte cage is
bicycle-free at radius $5$ at the same $\tfrac13$, and
generalized-polygon cages of degree $\bar d\sim s^{1/(R+1)}$ push the
density toward $1$ at fixed $R$. (Stochastic search over bicycle-free
graphs found exactly this ceiling, $0.344$ at $R=3$ in small sizes,
after first refuting the linear conjecture's banana extremal at
density $1/(R+1)$.) At the operating scale the lemma gives precisely
what is needed: for $s\le\ell$ and $R\ge C\log\ell/\delta$,
$\operatorname{rank}\le1+\delta\ell$.

\begin{proposition}[$\varepsilon$-theorem, bicycle-free regime]\label{thm:bf}
Fix $\delta\in(0,1)$, $L$, $\Lambda$, and let $G$ be $d$-regular on
$n$ vertices, bicycle-free at radius
$R\ge C\log\log n/\delta$, with $W$ spanned by any consistent family
of even cycles of length $\le L$ and $\Lambda_W\le\Lambda$. Then some
signing in the solution family satisfies
$\rho(A_\sigma)\le2\sqrt{d-1}\,(1+C\delta\log(1/\delta))(1+o(1))$.
\end{proposition}

\begin{proof}
Set $\ell=2\lceil\log n\rceil$ and $k=\lceil\log n\rceil$; we may
assume $R\le\ell$, since otherwise every support below is bicycle-free
at radius exceeding its size and Proposition~\ref{thm:eps} applies
with $r_0=1$. Fix $j\le2\ell$ and bound
$\sum_{z\in W}N^{\mathrm{nb}}_j(z)$ by rerunning
Propositions~\ref{prop:upper} and~\ref{prop:confined}. A counted walk
has support $S$ with $s\le(j+\Lambda)/2\le2\ell$ edges and minimum
degree $2$; by Lemma~\ref{lem:rank}, its cycle rank satisfies
$r\le1+Cs\log(2\ell)/R\le1+\delta(j+\Lambda)$ once
$R\ge C'\log\log n/\delta$. The fresh-run identity is untouched: the
walk performs exactly $r$ fresh runs, so the run-start schedule costs
$\binom{j}{r}\le(ej/r)^{r}\le(e/\delta)^{\delta(j+\Lambda)+1}$, and
the fresh choices cost $(d-1)^{s}\le(d-1)^{\Lambda/2}(d-1)^{j/2}$.
Each of the at most $r+1$ stale segments is a non-backtracking path in
$S$, and Lemma~\ref{lem:window} bounds its continuations by
$(2s(R+3))^{\lceil t_i/R\rceil}$ with $\sum_it_i\le j$, for a total
window exponent at most $j/R+\delta(j+\Lambda)+3$. Hence
\[
\sum_{z\in W}N^{\mathrm{nb}}_j(z)\;\le\;
K\,\hat\rho^{\,j},\qquad
\hat\rho:=(1+\eta)\sqrt{d-1},
\]
where $K\le nd\,(d-1)^{\Lambda/2}(e/\delta)^{\delta\Lambda+1}
\big(4\ell(R+3)\big)^{\delta\Lambda+3}$ collects all factors not
exponential in $j$, and
\[
\log(1+\eta)\;=\;\delta\log(e/\delta)
+\Big(\tfrac1R+\delta\Big)\log\!\big(4\ell(R+3)\big)
\;\le\;C\delta\log(1/\delta)
\]
at $R\ge C'\log\log n/\delta$ and $R\le\ell$. The peeling and transfer
of Lemma~\ref{lem:transfer} and the assembly in the proof of
Proposition~\ref{thm:eps} now apply verbatim with this $K$ and
$\hat\rho$, giving
$\mathbb E_{\mathcal F}\tr A_\sigma^{2k}\le
C''\,K\,(Ck)^{2}\big((1+\eta)^{2}\,2\sqrt{d-1}\big)^{2k}$, and
Corollary~\ref{cor:certificate} at $k=\lceil\log n\rceil$ yields
$\min_{\sigma\in\mathcal F}\rho(A_\sigma)\le
2\sqrt{d-1}\,(1+\eta)^{2}\big(C''K(Ck)^{2}\big)^{1/2k}$, which is the
claim.
\end{proof}

The extension from $z=0$ to all of $W$ costs only a constant. Set
$\Lambda_W:=\max_{z\in W}|z|\le L\dim W$.

\begin{proposition}[Confined-walk upper bound]\label{prop:confined}
Assume $\Lambda_W\le\ell$ and that $G$ is
$(r_0,\delta)$-subcritical at scale $\ell$. Then, with
$N^{\mathrm{nb}}_\ell(z)$ the number of closed non-backtracking walks
of length $\ell$ and parity $z$,
\[
\sum_{z\in W}N^{\mathrm{nb}}_\ell(z)\;\le\;
C_{r_0,\delta,\Lambda_W,d}\;n\,(C\ell)^{3(r_0+1)^2}
\big((1+\delta)^{2(r_0+2)}\sqrt{d-1}\big)^{\ell}.
\]
\end{proposition}

\begin{proof}
A counted walk has $|z(w)|\le\Lambda_W$; edges of $z(w)$ carry odd
multiplicity $\ge1$ and all others even multiplicity $\ge2$, so its
support has $s\le(\ell+\Lambda_W)/2$ edges and, as before, minimum
degree $2$. The proof of Proposition~\ref{prop:upper} used evenness
only through the bound on $s$; it applies verbatim, with
$(d-1)^{s}\le(d-1)^{\Lambda_W/2}(d-1)^{\ell/2}$ and the corresponding
constant absorbed into $C_{r_0,\delta,\Lambda_W,d}$.
\end{proof}

To reach the adjacency side, we use the classical linearization (see
e.g.\ \cite{bordenave20}). Define $\mathsf A_0=I$,
$\mathsf A_1=A_\sigma$, $\mathsf A_2=A_\sigma^2-dI$, and
$\mathsf A_{j+1}=A_\sigma\mathsf A_j-(d-1)\mathsf A_{j-1}$ for
$j\ge2$; then $(\mathsf A_j)_{uv}$ is the signed count of
non-backtracking walks $u\to v$ of length $j$, and inverting the
recursion,
$A_\sigma^{m}=\sum_{j\le m}\alpha_{m,j}\mathsf A_j$ where
$\alpha_{m,j}\ge0$ counts weight-$(d-1)$ downward-step walks on
$\mathbb N$ from $0$ to $j$. Both identities and the envelope below
were verified symbolically through $m=8$.

\begin{lemma}[Transfer envelope and peeling]\label{lem:transfer}
(i) $\sum_j\alpha_{m,j}(\sqrt{d-1})^{\,j}\le
\tfrac{d}{d-1}(m+1)\,\big(2\sqrt{d-1}\big)^{m}$.
(ii) Let $\widetilde N_j(z)$ count vertex-based non-backtracking closed
walks (backtracking permitted at the wrap) of length $j$ and parity
$z$. Then
$\widetilde N_j(z)\le\sum_{p\ge0}d(d-1)^{p}\,
N^{\mathrm{nb}}_{j-2p}(z)+d(d-1)^{j/2-1}\mathbf 1_{z=0}$.
\end{lemma}

\begin{proof}
(i) Each $\mathbb N$-walk of length $m$ ending at $j$ has $(m-j)/2$
down-steps, so $\alpha_{m,j}\le\binom{m}{(m+j)/2}(d-1)^{(m-j)/2}\cdot
\tfrac d{d-1}\le\tfrac d{d-1}2^m(d-1)^{(m-j)/2}$; multiply by
$(\sqrt{d-1})^j$ and sum. (ii) A vertex-based walk is either
cyclically non-backtracking or of the form $e\cdot w'\cdot\bar e$ with
$w'$ vertex-based of length $j-2$ at the head of $e$; peeling the
pendant edge adds $2$ to the multiplicity of $e$, leaving the parity
unchanged. Iterate: a pendant path of length $p$ costs at most
$d(d-1)^{p-1}$ choices, and the residue is cyclically non-backtracking
(or empty, forcing $z=0$).
\end{proof}

\begin{proposition}[$\varepsilon$-theorem, dilute regime]\label{thm:eps}
Fix $r_0,\delta,L,\Lambda$ and let $G$ be $d$-regular on $n$
vertices, $(r_0,\delta)$-subcritical at scale
$\ell=2\lceil\log n\rceil$.
Let $W$ be spanned by any consistent family of even cycles of length
$\le L$ with $\Lambda_W\le\Lambda$ --- possibly $W=\{0\}$. Then some
signing in the solution family $\mathcal F$ satisfies
\[
\rho(A_\sigma)\;\le\;
2\sqrt{d-1}\;(1+\delta)^{4(r_0+2)}
\Big(1+O_{r_0,\delta,\Lambda,d}\big(\tfrac{\log\log n}{\log n}\big)\Big).
\]
\end{proposition}

\begin{proof}
For $n$ large enough that $\Lambda\le\ell$,
Proposition~\ref{prop:confined} applies at every length $j\le2k$ by
monotonicity of the subcriticality notion. Write
$\hat\rho=(1+\delta)^{2(r_0+2)}\sqrt{d-1}$ and let $K$ denote
the prefactor of Proposition~\ref{prop:confined}. Since the recursion
defining $\mathsf A_j$ preserves the signed-walk interpretation, the
annihilator computation gives
$\mathbb E_{\mathrm{Ann}(W)}\tr\mathsf A_j(A_\sigma)
=\sum_{z\in W}\widetilde N_j(z)$. By
Lemma~\ref{lem:transfer}(ii), Proposition~\ref{prop:confined}, and
$\hat\rho^2\ge d-1$ (terms nonincreasing in $p$),
$\sum_{z\in W}\widetilde N_j(z)\le(j+2)\,dK\hat\rho^{\,j}$. Hence with
$k=\lceil\log n\rceil$,
\[
0\le\mathbb E_{\mathcal F}\tr A_\sigma^{2k}
\le\mathbb E_{\mathrm{Ann}(W)}\tr A_\sigma^{2k}
=\sum_j\alpha_{2k,j}\sum_{z\in W}\widetilde N_j(z)
\le (2k+2)dK\sum_j\alpha_{2k,j}\hat\rho^{\,j},
\]
and by Lemma~\ref{lem:transfer}(i) applied after extracting
$(1+\delta)^{2(r_0+2)j}\le(1+\delta)^{4(r_0+2)k}$, the right side is at
most $C'\,n\,(Ck)^{3(r_0+1)^2+2}
\big((1+\delta)^{4(r_0+2)}\,2\sqrt{d-1}\big)^{2k}$. Corollary
\ref{cor:certificate} and $\big(n(Ck)^{c}\big)^{1/2k}
=1+O(\log\log n/\log n)$ conclude.
\end{proof}

Taking $W=\{0\}$ recovers a near-Ramanujan statement for random
signings of dilute graphs in the spirit of \cite{mop20}, from entirely
different machinery; taking $W$ maximal uses the full parity family.
Random $2$-lift towers satisfy the hypotheses with $r_0=O(1)$,
$\delta=o(1)$, $\Lambda=O(1)$ with high probability, so the
$\varepsilon$-version of the Bilu--Linial conjecture holds along them.

\section{Outlook}

The target is now a single inequality: show
$R_{\mathcal F}(k)\le2\sqrt{d-1}+\varepsilon(d,\beta_L)$ at
$k\asymp\log n$, i.e.\ that the single-cycle classes
$-\sum_{C}N_{2k}(\chi_C)$ dominate the even-wrap excess
$N_{2k}(0)-n\,t_{2k}(d)$ while the pair classes stay controlled. The
structure of Proposition~\ref{prop:master} is that of an
Ihara--Bass/cluster expansion in which the local factors of the
constraint cycles appear with inverted sign; Proposition~\ref{prop:Lident}
makes this exact, and reduces the analytic core of the program to
Question~\ref{q:confined}, with the bicycle-free walk taxonomy of
\cite{mop20} supplying the combinatorics
and now entering with cancellation rather than by union bound. The
variance-onset data locate the error term: long-cycle fluctuations of
standard deviation $\approx10^{-2}$, persistent in $n$, which the
$k\asymp\log n$ analysis must absorb.

\begin{thebibliography}{9}
\bibitem{ahl02} N.~Alon, S.~Hoory, N.~Linial, \emph{The Moore bound
for irregular graphs}, Graphs Combin.~18 (2002), 53--57.
\bibitem{bass92} H.~Bass, \emph{The Ihara--Selberg zeta function of a
tree lattice}, Internat.~J.~Math.~3 (1992), 717--797.
\bibitem{bilulinial06} Y.~Bilu, N.~Linial,
\emph{Lifts, discrepancy and nearly optimal spectral gap},
Combinatorica 26 (2006), 495--519.
\bibitem{bordenave20} C.~Bordenave, \emph{A new proof of Friedman's
second eigenvalue theorem and its extension to random lifts},
Ann.~Sci.~\'Ec.~Norm.~Sup\'er.~53 (2020), 1393--1439.
\bibitem{godsilgutman81} C.~D.~Godsil, I.~Gutman,
\emph{On the matching polynomial of a graph}, in: Algebraic Methods in
Graph Theory, 1981.
\bibitem{hastad01} J.~H\aa stad, \emph{Some optimal inapproximability
results}, J.~ACM 48 (2001), 798--859.
\bibitem{heilmannlieb72} O.~J.~Heilmann, E.~H.~Lieb,
\emph{Theory of monomer-dimer systems}, Comm.~Math.~Phys.~25 (1972),
190--232.
\bibitem{huang19} H.~Huang, \emph{Induced subgraphs of hypercubes and a
proof of the Sensitivity Conjecture}, Ann.~of Math.~190 (2019), 949--955.
\bibitem{lieb94} E.~H.~Lieb, \emph{Flux phase of the half-filled band},
Phys.~Rev.~Lett.~73 (1994), 2158--2161.
\bibitem{mss15} A.~W.~Marcus, D.~A.~Spielman, N.~Srivastava,
\emph{Interlacing families I: Bipartite Ramanujan graphs of all degrees},
Ann.~of Math.~182 (2015), 307--325.
\bibitem{mop20} S.~Mohanty, R.~O'Donnell, P.~Paredes,
\emph{Explicit near-Ramanujan graphs of every degree}, STOC 2020;
arXiv:1909.06988.
\bibitem{xuzhang26} Z.~Xu, X.~Zhang, \emph{An improved upper bound for
the Bilu--Linial conjecture}, arXiv:2606.28797 (2026).
\end{thebibliography}

\end{document}
